% META: {"id": "TSPE-COMBI-EX-013", "chapitre": "TSPE-COMBINATOIRE", "type_objet": "exercice", "capacites": ["TSPE-COMBINATOIRE-C3"], "capacites_codes": ["C3"], "methodes": ["M1"], "parcours": 1, "competences": ["calculer"], "duree_min": 8, "mode_creation": "ex_nihilo", "sources_inspiration": [], "fichier_tex": "chapitres/TSPE-COMBINATOIRE/exercices/TSPE-COMBI-EX-013.tex", "corrige_tex": "chapitres/TSPE-COMBINATOIRE/corriges/TSPE-COMBI-CO-013.tex", "status": "approved"}
% BEGIN-VERIFY
% from sympy import binomial
% assert binomial(6, 2) + binomial(6, 3) == binomial(7, 3) == 35
% END-VERIFY
\begin{exercice}{TSPE-COMBI-EX-006}{2}{10}
Sachant que $\dbinom{6}{2}=15$ et $\dbinom{6}{3}=20$, calculer $\dbinom{7}{3}$ sans utiliser la formule factorielle, en utilisant la relation de Pascal.
\end{exercice}
